\subsection{Controlled Severity-Score and State-Carry Study}
\label{sec:component-study}

The continuous comparison evaluates complete procedures. To distinguish
two mechanisms within the physical bank, we fixed a subsequent $2\times2$
study: marginal amplitude evidence versus a bounded profile score, and
physical-state carry versus removal at inferred closure. The five network
models, observations, startup fits, operators, proposal rules and
publication thresholds are common to all four arms. Identical thresholds
do not imply identical false-alarm rates; both coverage and alarm counts
must be measured.

For one load history and amplitude interval $[l,u]$, write its accumulated
conditional log likelihood as
\begin{equation}
 \ell_k(a)=c_k+b_k a-\tfrac12 G_k a^2.
 \label{eq:ablation-quadratic}
\end{equation}
The original score uses $B_k=\log\int_l^u
\exp(\ell_k(a)-c_k)p(a)\,da$ with the existing truncated Gaussian prior.
The alternative uses
\begin{equation}
 \widehat a_k=\operatorname{clip}_{[l,u]}(b_k/G_k),\qquad
 Q_k=b_k\widehat a_k-\tfrac12G_k\widehat a_k^2,
 \label{eq:ablation-profile}
\end{equation}
with the constant-likelihood case treated separately. Each report adds
$B_k-B_{k-1}$ or $Q_k-Q_{k-1}$, so a rolling replay does not repeatedly
add the same suffix evidence. The conditional Gaussian amplitude/state
updates remain unchanged. Profile weights are normalized scores, not
posterior probabilities. This intervention tests the amplitude-scoring
choice, rather than removing every Bayesian calculation.
The existing Gaussian sufficient statistics supply
$G_k=v_k^{-1}-v_0^{-1}$ and $b_k=\mu_k/v_k$, removing the computational
prior precision before profiling. Fixed line-opening amplitudes and the
prior weights of amplitude intervals are unchanged.

At an internally inferred closure, the carry arm retains the projected
augmented state. The reset arm sets its 114 estimated physical-state
deviations to zero and retains instrumental/discrepancy coordinates,
covariance regime and network-model weights. This is an estimator
ablation: every observed record comes from the same continuous physical
trajectory. It does not test removal of all temporal information.

The protocol uses all 18 matched-network continuous parents at the
original noise level, plus the first three normal and persistent-line
parents by identifier. It was fixed before ablation outcomes, but the
underlying records had already been examined. Results are therefore a
posthoc mechanism study, with paired intervals over physical parents,
not a new independent confirmation.
